% ID: FIS-FLU-PID-001
% Titolo: Pressione idrostatica in due recipienti cilindrici
% Disciplina: Fisica
% Area: Fluidostatica
% Argomento: Pressione idrostatica
% Sottoargomento: Pressione sul fondo di un recipiente
% Classe: Seconda liceo scientifico
% Difficolta: 3
% Tipo: problema_numerico
% Risultato: soluzione da aggiungere
% Tempo_stimato: 8 min
% Competenze: modellizzazione, proporzionalita, uso delle unita di misura
% Prerequisiti: densita, pressione idrostatica
% Tag: fisica, fluidostatica, pressione_idrostatica, acqua, recipienti
% Fonte: verifica_fisica_fluidostatica_anonimizzata
% Autore: Riccardo Carli
% Licenza: CC-BY-SA-4.0

\begin{esercizio}
Due recipienti cilindrici completamente pieni d'acqua, di densita
$\rho = 1000\,\mathrm{kg/m^3}$, hanno raggio di base rispettivamente
$4\,\mathrm{cm}$ e $6\,\mathrm{cm}$; l'altezza del primo recipiente e
$20\,\mathrm{cm}$.

Calcola la pressione idrostatica esercitata sul fondo del primo recipiente e
l'altezza del secondo recipiente, sapendo che la pressione esercitata sul suo
fondo e il doppio di quella esercitata sul fondo del primo.
\end{esercizio}

\begin{soluzione}
% Soluzione da aggiungere.
\end{soluzione}

